% brief/02_physics.tex: the data in brief, then why the hand is detectable. Wording from sections/02_data.tex
% and sections/03_physics.tex, shortened; PD01 is left out of the brief. The video footnote is the one of
% sections/02_data.tex. Figure 2 (both columns, natural 6.2-in width) is defined before \section so that it is met on
% page 1 and lands at the top of page 2 (brief.tex: one figure* per page top).
\begin{figure*}[tbp]
\centering
\genfig{core/S2_13_first_half_yaw}
\caption{The yaw about gravity accumulated over the first half of the approach separates the hands (left: AUC within
each person; right: the pooled distributions). The net yaw over the whole approach (grey) has a lower pooled AUC,
\R{describe/net_yaw_auc_pooled}, not tested in a paired comparison.}
\label{fig:S2_13_first_half_yaw}
\end{figure*}

\section{Why the hand is detectable}\label{sec:physics}
The data come from \R{describe/n_persons} students (\R{describe/n_left_handed} left-handed;
\R{describe/persons_android} on Android, \R{describe/persons_ios} on iOS), who followed the prompts of the NoseTouch
web app for the hand, the grip (upright or sideways) and the speed of each trial.\footnote{\raggedright A video of a
data-collection session: \url{https://drive.google.com/file/d/1V9PF623HkJ8MjR6loJhUHnDkCx60W-Jl/view?usp=sharing}.}
A pre-specified exclusion cascade keeps
\R{describe/primary_trials} instructed trials and \R{describe/natural_use_n} free-choice ones of the
\R{describe/attempts_total} recorded attempts.

In the free-choice trial each person chose the hand and reported it afterwards; these natural-use trials, all
right-hand, are scored but never used for training (full report, Section~\fullref{sec:data}).

Mirroring the body across its midline swaps the hands. It flips the sign of chiral (mirror-odd) quantities, such as
the yaw rate about gravity, and leaves mirror-even ones, such as speed and duration, unchanged.

If the two hands move as mirror images, only chiral quantities can carry the hand. The yaw rate about gravity is also
unchanged by any grip, because a grip only rotates the phone's axes.

Early in the approach the forearm swings toward the body's midline, so, seen from above, the phone turns
counter-clockwise in the right hand and clockwise in the left.

The yaw about gravity accumulated over the first half of the approach separates the hands with a pooled area under
the ROC curve (AUC) of \R{describe/fhy_auc_pooled}, and within persons the AUC is above 0.5 for
\R{describe/fhy_persons_auc_above_half} of \R{describe/n_persons} persons (Figure~\ref{fig:S2_13_first_half_yaw}).
Its sign alone, with no training, classifies \R{describe/fhy_sign_rule_accuracy} of the trials.

The models therefore get time-resolved channels in a gravity frame, and a body frame that follows the phone's
heading adds channels that do not depend on an upright or sideways grip (a screen-away hold needs a separate check,
Section~\ref{sec:results}).
